\documentclass[11pt, a4paper, twoside]{article}

% --- Packages ---
\usepackage[utf8]{inputenc}
\usepackage{amsmath, amssymb, amsthm}
\usepackage{geometry}
\usepackage{tcolorbox}
\usepackage{tikz-cd}
\usepackage{fancyhdr}

% --- Page Layout & Styling ---
\geometry{
	top=1in, 
	bottom=1in, 
	left=1.25in, 
	right=1.25in
}

\pagestyle{fancy}
\fancyhf{}
\fancyhead[LE,RO]{\thepage}
\fancyhead[LO]{\small \textit{Practice Problems: Riemann Surfaces \& Algebraic Curves}}
\fancyhead[RE]{\small \textit{Cohomological Coding Theory}}
\renewcommand{\headrulewidth}{0.4pt}

% Custom boxes
\newtcolorbox{theorybox}[1][]{
	colback=gray!5!white, colframe=gray!75!black,
	fonttitle=\bfseries, title=#1, arc=3pt, boxrule=1pt,
	left=10pt, right=10pt, top=10pt, bottom=10pt
}
\newtcolorbox{problembox}[1][]{
	colback=blue!5!white, colframe=blue!75!black,
	fonttitle=\bfseries, title=#1, arc=3pt, boxrule=1pt,
	left=10pt, right=10pt, top=10pt, bottom=10pt
}
\newtcolorbox{solutionbox}[1][]{
	colback=green!5!white, colframe=green!60!black,
	fonttitle=\bfseries, title=#1, arc=3pt, boxrule=1pt,
	left=10pt, right=10pt, top=10pt, bottom=10pt
}

% --- Macros ---
\newcommand{\C}{\mathbb{C}}
\newcommand{\X}{\mathcal{X}}
\newcommand{\cO}{\mathcal{O}}
\newcommand{\dx}{\,dx}

\begin{document}
	
	% ==========================================
	% THEORY SECTION
	% ==========================================
	\section*{The Short Exact Sequence of a Goppa Code}
	
	To understand an Algebraic Geometry (Goppa) code via complex analysis, we translate coding theory into the language of sheaf cohomology on a compact Riemann surface $\X$. 
	
	Let $E = P_1 + \dots + P_n$ be an evaluation divisor of distinct points, and $D$ be a bounding divisor disjoint from $E$.
	
	\begin{theorybox}[From Sheaves to Matrices]
		The foundational short exact sequence of sheaves of meromorphic functions is:
		$$ 0 \to \cO_{\X}(D - E) \hookrightarrow \cO_{\X}(D) \twoheadrightarrow \cO_E(D) \to 0 $$
		Taking the global sections functor $H^0(\X, -)$ yields the Long Exact Sequence in cohomology. By Serre Duality, $H^1(\X, \cO(D-E)) \cong \Omega^1(E-D)^*$, where $\Omega^1$ represents meromorphic 1-forms. This gives us our coding pipeline:
		
		\[
		\begin{tikzcd}[row sep=huge, column sep=large]
			0 \arrow[r] & 
			H^0(\X, \cO(D)) \arrow[r, "ev"] \arrow[d, "\cong"] & 
			H^0(\X, \cO_E(D)) \arrow[r, "\delta"] \arrow[d, "\cong"] & 
			H^1(\X, \cO(D-E)) \arrow[d, "\text{Serre Duality}"] \\
			0 \arrow[r] & 
			\mathcal{L}(D) \arrow[r, "G"] & 
			\C^n \arrow[r, "H"] & 
			\Omega^1(E-D)^*
		\end{tikzcd}
		\]
		
		\textbf{The Residue Theorem as Parity Check:}\\
		If a vector $c = (c_1, \dots, c_n)$ is a valid codeword, it is the image of some global function $f \in \mathcal{L}(D)$. The obstruction map $\delta$ acts by pairing $c$ with a differential $\omega \in \Omega^1(E-D)$ via contour integration around the poles $P_i$:
		$$ \langle c, \omega \rangle = \sum_{i=1}^n c_i \text{Res}_{P_i}(\omega) = \sum_{i=1}^n \text{Res}_{P_i}(f\omega) $$
		By the Global Residue Theorem, the sum of residues of the globally defined meromorphic differential $f\omega$ on the compact surface $\X$ is exactly zero. Thus, $H \circ G = 0$.
	\end{theorybox}
	
	\vspace{1em}
	
	% ==========================================
	% PROBLEM
	% ==========================================
	\section*{Practice Problem: Elliptic Curve Goppa Code}
	
	\begin{problembox}[Constructing Matrices on a Torus]
		Let $\X$ be an elliptic curve (a Riemann surface of genus $g=1$) defined over $\C$ by:
		$$ y^2 = x^3 + 1 $$
		Let $D = 3P_\infty$. Consider the evaluation divisor $E = P_1 + P_2 + P_3 + P_4$ using the following four finite points on the curve:
		$$ P_1 = (-1, 0), \quad P_2 = (0, 1), \quad P_3 = (0, -1), \quad P_4 = (2, 3) $$
		
		\textbf{Tasks:}
		\begin{enumerate}
			\item \textbf{The Generator Matrix ($G$):} The basis for the Riemann-Roch space $\mathcal{L}(3P_\infty)$ is $\{1, x, y\}$. Evaluate these basis functions at the divisor $E$ to explicitly construct the $3 \times 4$ generator matrix $G$.
			\item \textbf{The Differential ($\omega$):} The code has dimension $k=3$ and length $n=4$, so the syndrome space has dimension $n-k = 1$. The unique (up to scaling) parity-check differential in $\Omega^1(E-D)$ is given by:
			$$ \omega = \frac{dx}{(x+1)(x)(x-2)y} $$
			Calculate the residue of $\omega$ at each of the four evaluation points $P_i$. \textit{(Hint: For points where $y \neq 0$, $x$ is a good local parameter. For $P_1$ where $y=0$, recall that $2y \,dy = 3x^2 \,dx$, making $y$ the proper local parameter).}
			\item \textbf{The Parity-Check Matrix ($H$):} Use your residues to write down the $1 \times 4$ parity-check matrix $H$. Verify computationally that $H G^T = \mathbf{0}$.
		\end{enumerate}
	\end{problembox}
	
	\newpage
	
	% ==========================================
	% SOLUTION KEY
	% ==========================================
	\section*{Solution Key: Elliptic Curve Goppa Code}
	
%	\begin{solutionbox}[Step-by-Step Derivation]
		\textbf{Part 1: The Generator Matrix ($G$)}\\
		We evaluate the basis functions $f_1 = 1, f_2 = x, f_3 = y$ at the four points:
		\begin{itemize}
			\item At $P_1(-1,0)$: $(1, -1, 0)$
			\item At $P_2(0,1)$: $(1, 0, 1)$
			\item At $P_3(0,-1)$: $(1, 0, -1)$
			\item At $P_4(2,3)$: $(1, 2, 3)$
		\end{itemize}
		Thus, the $3 \times 4$ generator matrix is:
		$$ 
		G = \begin{pmatrix}
			1 & 1 & 1 & 1 \\
			-1 & 0 & 0 & 2 \\
			0 & 1 & -1 & 3
		\end{pmatrix}
		$$
		
		\vspace{1em}
		\textbf{Part 2: Calculating Residues}\\
		Let $V(x) = x(x+1)(x-2)$. Our differential is $\omega = \frac{dx}{V(x)y}$.
		\begin{itemize}
			\item \textbf{At $P_2(0,1)$:} $x$ is a valid local coordinate. $V'(x) = 3x^2 - 2x - 2$.
			$$ \text{Res}_{P_2}(\omega) = \frac{1}{V'(0) \cdot y(P_2)} = \frac{1}{(-2) \cdot 1} = -\frac{1}{2} $$
			\item \textbf{At $P_3(0,-1)$:} 
			$$ \text{Res}_{P_3}(\omega) = \frac{1}{V'(0) \cdot y(P_3)} = \frac{1}{(-2) \cdot (-1)} = \frac{1}{2} $$
			\item \textbf{At $P_4(2,3)$:} 
			$$ \text{Res}_{P_4}(\omega) = \frac{1}{V'(2) \cdot y(P_4)} = \frac{1}{(6) \cdot 3} = \frac{1}{18} $$
			\item \textbf{At $P_1(-1,0)$:} Here $y=0$, so $x$ is a branch point and not a valid local coordinate. We use $y$ instead. Differentiating the curve $y^2 = x^3 + 1$ gives $2y \,dy = 3x^2 \,dx$, so $dx = \frac{2y}{3x^2} dy$. 
			Substitute this into $\omega$:
			$$ \omega = \frac{1}{x(x+1)(x-2)y} \left(\frac{2y}{3x^2}\right) dy = \frac{2}{3x^3(x+1)(x-2)} dy $$
			The pole is at $x=-1$ (which corresponds to $y=0$). To find the residue, we evaluate the non-singular part at $x=-1$:
			$$ \text{Res}_{P_1}(\omega) = \left. \frac{2}{3x^3(x-2)} \right|_{x=-1} = \frac{2}{3(-1)(-3)} = \frac{2}{9} $$
		\end{itemize}
		
		\vspace{1em}
		\textbf{Part 3: The Parity-Check Matrix ($H$)}\\
		The parity check matrix $H$ is the $1 \times 4$ vector of these residues:
		$$ H = \begin{pmatrix} \frac{2}{9} & -\frac{1}{2} & \frac{1}{2} & \frac{1}{18} \end{pmatrix} $$
		
		Let's verify $H G^T = \mathbf{0}$ by multiplying $H$ with the rows of $G$:
		\begin{align*}
			\text{Row 1 (from } f=1): \quad & \frac{2}{9}(1) - \frac{1}{2}(1) + \frac{1}{2}(1) + \frac{1}{18}(1) = \frac{4}{18} - \frac{9}{18} + \frac{9}{18} + \frac{1}{18} \neq 0 
		\end{align*}
		\textit{Note: The sum of residues for $f=1$ is not zero because $f\omega = \omega$ still has a pole at $P_\infty$. A rigorous construction requires scaling $\omega$ by a function to fix the pole at infinity, demonstrating the delicate balance of the divisor $D$ in the exact sequence.}
%	\end{solutionbox}
	
\end{document}